%% ma: L10-B02-0045
%% lop: 10
%% chuong: 1
%% bai: 2
%% dang: 10.2.2
%% loai: TN4
%% muc: TH
%% dap_an: "-"
%% nguon: "Ảnh giáo viên gửi 10/10/2026 (ThS. Nguyễn Hoàng Việt, Chương 2 Tập hợp), Đề 2 (đề thứ hai, không ghi tiêu đề), Phần 1 Câu 2"
%% kien_thuc: "Bài 2 (tập hợp và các phép toán trên tập hợp)"
%% trang_thai: cach-ly
%% ly_do: "Theo đề, X = {1; 2} và S = 3 (kiểm bằng sympy), không có trong bốn phương án 9/2, 5, 6, 4. Gợi ý: nếu thừa số thứ ba là 2x^2 - 7x + 3 (nghiệm 1/2 và 3) thì X = {1; 2; 3} và S = 6 (phương án C); đối chiếu sách gốc"
\begin{ex}
Cho tập $X=\{x\in\mathbb N\mid(x^2-4)(x-1)(x^2-7x+3)=0\}$. Tính tổng $S$ các phần tử của $X$.
\choice{$S=\dfrac92$}{$S=5$}{$S=6$}{$S=4$}
\loigiai{$x^2-4=0\Leftrightarrow x=\pm2$; $x-1=0\Leftrightarrow x=1$; $x^2-7x+3=0$ có nghiệm $\dfrac{7\pm\sqrt{37}}{2}$ không nguyên. Nghiệm tự nhiên: $X=\{1;2\}$, $S=3$ (không có trong các phương án).}
\end{ex}
